\subsection{可表现环}

称环 A 为可表现的，如果存在一个正则环 R 以及一个从 R 到 A 的满射同态。

按定义，正则环都是可表现的。

命题 6.--- a) 可表现环的每个分式环都是可表现的。可表现环上的每个有限生成代数都是可表现环。

\begin{enumerate}
\def\labelenumi{\alph{enumi})}
\setcounter{enumi}{1}
\item
  设 A 为可表现环，J 为 A 的理想。A 关于 J-进拓扑的分离完备化 \(\widehat{\mathbf{A}}\) 是可表现的。
\item
  每个完备 Noether 局部环都是可表现的。
\item
  设 A 为可表现局部环。存在一个正则局部环 R 以及一个从 R 到 A 的满射局部同态。
\end{enumerate}

设 A 为可表现环；选取一个正则环 R 以及一个满射同态 \(\rho: R \to A\) 。

\begin{enumerate}
\def\labelenumi{\alph{enumi})}
\tightlist
\item
  设 S 为 A 的一个乘性子集；令 \(T = \rho^{-1}(S)\) 。同态 \(T^{-1}R \to S^{-1}A\)（由 \(\rho\) 导出）是满射的，且环 \(T^{-1}R\) 是正则的 (第 3 节, 例 1)，因此 \(S^{-1}A\) 是可表现的。
\end{enumerate}

设 B 为有限生成 A-代数；存在一个有限集 I 以及一个满射同态 \(\mathrm{A}[(\mathrm{T}_{i})_{i\in\mathbb{I}}]\to\mathrm{B}\) ，从而有一个满射同态 \(\mathrm{R}[(\mathrm{T}_{i})_{i\in\mathbb{I}}]\to\mathrm{B}\) 。由于环 \(\mathrm{R}[(\mathrm{T}_{i})_{i\in\mathbb{I}}]\) 是正则的 (第 2 节, 命题 4 的推论 5)，环 B 是可表现的。

\begin{enumerate}
\def\labelenumi{\alph{enumi})}
\setcounter{enumi}{1}
\item
  令 \(1 = \rho^{-1}(\mathrm{J})\)，并用 \(\widehat{\mathbb{R}}\) 表示 \(\mathbb{R}\) 关于 I-进拓扑的分离完备化；对每个整数 \(n \geqslant 0\)，典范映射 \(\mathrm{I}^n / \mathrm{I}^{n+1} \to \mathrm{J}^n / \mathrm{J}^{n+1}\) 是满射的。于是，同态 \(\widehat{\mathbb{R}} \to \widehat{\mathbb{A}}\)（由 \(\rho\) 导出）是满射的 (III, § 2, 第 8 节, 定理 1 的推论 2)；由于 \(\widehat{\mathbb{R}}\) 正则 (第 2 节, 命题 4 的推论 3)，环 \(\widehat{\mathbb{A}}\) 是可表现的。
\item
  这由 IX, § 2, 第 5 节, 定理 3 a) 及 IX, § 3, 第 3 节, 定理 2 a) 得出。
\item
  设 \(\mathfrak{m}_{\Lambda}\) 为 A 的极大理想；则 \(\mathfrak{p} = \mathfrak{\rho}^{-1}(\mathfrak{m}_{\Lambda})\) 是 R 的素理想，局部环 \(R_{p}\) 正则，且同态 \(R_{p} \to A\)（由 \(\mathfrak{p}\) 导出）是局部的且是满射的。
\end{enumerate}

由于域和 Dedekind 环都是正则的，从而是可表现的，命题 6 蕴含代数几何中常见的大多数环都是可表现的。

命题 7. - 设 A 为可表现环。

\begin{enumerate}
\def\labelenumi{\alph{enumi})}
\item
  环 Λ 是 Noether 的且是级链的。
\item
  设 M 为有限生成 A-模。映射

\[
\mathfrak {p} \mapsto \dim_ {\mathrm{A} _ {\mathfrak {p}}} (\mathrm{M} _ {\mathfrak {p}}) - \operatorname{prof} _ {\mathrm{A} _ {\mathfrak {p}}} (\mathrm{M} _ {\mathfrak {p}})
\]

从 Spec(A) 到 Z 是上半连续的。

\item
  设 M 为 \(\Lambda\) 上的有限生成模。由 \(\mathfrak{p} \in \operatorname{Spec}(\Lambda)\)（使 \(\mathrm{A}_{\mathfrak{p}}\) -模 \(\mathrm{M}_{\mathfrak{p}}\) 是 Macaulay 的）组成的集合是一个稠密开集。它与 \(\operatorname{Supp}(\mathrm{M})\) 的交在 \(\operatorname{Supp}(\mathrm{M})\) 中稠密。
\end{enumerate}

我们选取一个正则环 R 以及一个满射同态 R → A。

\begin{enumerate}
\def\labelenumi{\alph{enumi})}
\item
  环 R 是 Macaulay 的 (第 2 节, 命题 4 的推论 2)，因此 A 是级链的 (§ 2, 第 2 节, 命题 2 及 VIII, § 1, 第 3 节, 注 2)。
\item
  R-模 M 是有限生成的且具有有限投射维数 \(< +\infty\) (第 1 节, 命题 1)。把 \(\operatorname{Spec}(A)\) 等同于 \(\operatorname{Spec}(R)\) 的一个闭子集。于是函数 \(\mathfrak{p} \mapsto \dim_{A_{\mathfrak{p}}} (M_{\mathfrak{p}}) - \operatorname{prof}_{A_{\mathfrak{p}}} (M_{\mathfrak{p}})\) 在 \(\operatorname{Spec}(A)\) 上是函数 \(\mathfrak{q} \mapsto \dim_{R_{\mathfrak{q}}} (M_{\mathfrak{q}}) - \operatorname{prof}_{R_{\mathfrak{q}}} (M_{\mathfrak{q}})\) 在 \(\operatorname{Spec}(R)\) 上的限制。于是只需应用 § 3, 第 5 节的定理 1 的推论 4。
\item
  这按同上引文 c) 的方式证明。
\end{enumerate}

